\subsection{A.2 LSTM Implementation}
 
\noindent \textbf{Πηγή:} \texttt{Protein\_Network\_FIXED.ipynb}
 
 
\subsubsection*{Temporal Protein Network - FIXED VERSION}
 
\subsubsection*{1. Setup}
 
\begin{lstlisting}
import numpy as np
import pandas as pd
import torch
import torch.nn as nn
import torch.nn.functional as F
from torch.utils.data import Dataset, DataLoader
from sklearn.model_selection import StratifiedKFold
from sklearn.preprocessing import StandardScaler, LabelEncoder
from sklearn.feature_selection import SelectKBest, f_classif
from sklearn.impute import SimpleImputer
from sklearn.metrics import classification_report, confusion_matrix, accuracy_score
import matplotlib.pyplot as plt
import seaborn as sns
import warnings
warnings.filterwarnings('ignore')
 
# Device
device = torch.device('cuda' if torch.cuda.is_available() else 'cpu')
print(f"Using device: {device}")
 
# Reproducibility
def set_seed(seed=42):
    np.random.seed(seed)
    torch.manual_seed(seed)
    if torch.cuda.is_available():
        torch.cuda.manual_seed_all(seed)
 
set_seed(42)
print("Setup complete ✓")
\end{lstlisting}
 
\subsubsection*{2. Configuration}
 
\begin{lstlisting}
class Config:
    RANDOM_SEED = 42
    N_FOLDS = 5
    N_TOP_FEATURES = 50
 
    # Model - Small
    HIDDEN_DIM = 32
    DROPOUT = 0.5
 
    # Training - Conservative
    BATCH_SIZE = 8
    LEARNING_RATE = 1e-4
    WEIGHT_DECAY = 1e-3
    EPOCHS = 50
    PATIENCE = 10
 
    DEVICE = device
 
print("Config loaded ✓")
\end{lstlisting}
 
\subsubsection*{3. Data Loading}
 
\begin{lstlisting}
# ============================================================
# ΑΛΛΑΞΕ ΑΥΤΑ ΤΑ PATHS ΣΤΑ ΔΙΚΑ ΣΟΥ ΑΡΧΕΙΑ
# ============================================================
SAMPLES_PATH = '/content/DIOPTRA_Samples_list_Oncology.xlsx'
PROTEINS_PATH = '/content/DIOPTRA_Protein_results_Oncology.xlsx'
# ============================================================
\end{lstlisting}
 
\begin{lstlisting}
def load_data(samples_path, proteins_path, include_t00=True):
    """
    Φορτώνει τα δεδομένα.
    """
    samples = pd.read_excel(samples_path)
    proteins = pd.read_excel(proteins_path)
    data = samples.merge(proteins, on='SampleID')
 
    protein_cols = [c for c in proteins.columns if c != 'SampleID']
 
    # Paired samples (T0 → T1)
    t0_data = data[data['Timepoint'] == 'T0'].copy()
    t1_data = data[data['Timepoint'] == 'T1'].copy()
 
    t0_data['PatientID'] = t0_data['SampleID']
    t1_data['PatientID'] = t1_data['SampleID'].str.replace('-F', '', regex=False)
 
    paired = t0_data.merge(
        t1_data[['PatientID'] + protein_cols],
        on='PatientID',
        suffixes=('_T0', '_T1')
    )
    print(f"Paired samples: {len(paired)}")
 
    # Add T00 samples
    if include_t00:
        t00_data = data[data['Timepoint'] == 'T00'].copy()
        for col in protein_cols:
            t00_data[f'{col}_T0'] = t00_data[col]
            t00_data[f'{col}_T1'] = t00_data[col]
        t00_data['PatientID'] = t00_data['SampleID']
 
        cols_needed = ['PatientID', 'Group'] + [f'{c}_T0' for c in protein_cols] + [f'{c}_T1' for c in protein_cols]
        combined = pd.concat([paired[cols_needed], t00_data[cols_needed]], ignore_index=True)
        print(f"After adding T00: {len(combined)}")
    else:
        combined = paired
 
    # Remove NAA
    combined = combined[combined['Group'] != 'NAA']
    print(f"After removing NAA: {len(combined)}")
 
    # Create arrays
    X_T0 = combined[[f'{c}_T0' for c in protein_cols]].values
    X_T1 = combined[[f'{c}_T1' for c in protein_cols]].values
    X = np.stack([X_T0, X_T1], axis=1)  # (n_samples, 2, n_proteins)
 
    # Labels
    label_encoder = LabelEncoder()
    y = label_encoder.fit_transform(combined['Group'].values)
 
    print(f"\nClass distribution:")
    for i, cls in enumerate(label_encoder.classes_):
        print(f"  {cls}: {(y == i).sum()}")
 
    return X, y, protein_cols, label_encoder
 
 
# Load data
X_raw, y, protein_names, label_encoder = load_data(SAMPLES_PATH, PROTEINS_PATH)
class_names = list(label_encoder.classes_)
 
print(f"\nX shape: {X_raw.shape}")
print(f"Classes: {class_names}")
\end{lstlisting}
 
\subsubsection*{4. Handle NaN Values}
 
\begin{lstlisting}
# Check for NaN
print(f"NaN values in X_raw: {np.isnan(X_raw).sum()}")
print(f"Shape before: {X_raw.shape}")
 
n_samples, n_timepoints, n_proteins = X_raw.shape
 
# Step 1: Find proteins that have at least some valid values
valid_protein_mask = ~np.all(np.isnan(X_raw), axis=(0, 1))  # True if protein has any valid value
print(f"\nProteins with ALL NaN (will be removed): {(~valid_protein_mask).sum()}")
print(f"Valid proteins: {valid_protein_mask.sum()}")
 
# Step 2: Keep only valid proteins
X_raw = X_raw[:, :, valid_protein_mask]
protein_names = [p for p, valid in zip(protein_names, valid_protein_mask) if valid]
 
print(f"Shape after removing all-NaN proteins: {X_raw.shape}")
 
# Step 3: Impute remaining NaN with median
n_samples, n_timepoints, n_proteins = X_raw.shape
X_imputed = np.zeros_like(X_raw)
 
imputer = SimpleImputer(strategy='median')
 
for t in range(n_timepoints):
    X_imputed[:, t, :] = imputer.fit_transform(X_raw[:, t, :])
 
X_raw = X_imputed
 
print(f"\nNaN values after imputation: {np.isnan(X_raw).sum()}")
print(f"Final shape: {X_raw.shape}")
\end{lstlisting}
 
\subsubsection*{5. Feature Selection}
 
\begin{lstlisting}
def select_top_features(X, y, protein_names, n_features=50):
    """
    Επιλέγει τις top-K πρωτεΐνες με βάση ANOVA F-score.
    """
    # Use T0 for feature selection
    X_t0 = X[:, 0, :]
 
    # Adjust n_features if we have fewer proteins
    n_features = min(n_features, X_t0.shape[1])
 
    # Select top features
    selector = SelectKBest(f_classif, k=n_features)
    selector.fit(X_t0, y)
 
    # Get selected indices
    selected_mask = selector.get_support()
    selected_indices = np.where(selected_mask)[0]
 
    # Apply to both timepoints
    X_selected = X[:, :, selected_indices]
 
    # Get selected protein names
    selected_proteins = [protein_names[i] for i in selected_indices]
 
    return X_selected, selected_indices, selected_proteins, selector
 
 
# Apply feature selection
X, selected_indices, selected_proteins, feature_selector = select_top_features(
    X_raw, y, protein_names, n_features=Config.N_TOP_FEATURES
)
 
print(f"Selected {len(selected_proteins)} proteins")
print(f"New X shape: {X.shape}")
print(f"\nTop 10 selected proteins:")
for i, p in enumerate(selected_proteins[:10]):
    print(f"  {i+1}. {p}")
\end{lstlisting}
 
\subsubsection*{6. Normalize Data}
 
\begin{lstlisting}
def normalize_data(X):
    """Normalize per feature across all samples and timepoints"""
    n_samples, n_timepoints, n_features = X.shape
    X_flat = X.reshape(-1, n_features)
 
    scaler = StandardScaler()
    X_flat_scaled = scaler.fit_transform(X_flat)
 
    # Clip extreme values
    X_flat_scaled = np.clip(X_flat_scaled, -5, 5)
 
    return X_flat_scaled.reshape(n_samples, n_timepoints, n_features), scaler
 
 
X, scaler = normalize_data(X)
print(f"Data normalized.")
print(f"Range: [{X.min():.2f}, {X.max():.2f}]")
print(f"Any NaN: {np.isnan(X).any()}")
\end{lstlisting}
 
\subsubsection*{7. Dataset Class}
 
\begin{lstlisting}
class ProteinDataset(Dataset):
    def __init__(self, X, y, mode='temporal'):
        self.y = torch.LongTensor(y)
 
        if mode == 'temporal':
            self.X = torch.FloatTensor(X)
        elif mode == 'static':
            self.X = torch.FloatTensor(X[:, 0:1, :])
        elif mode == 'delta':
            delta = X[:, 1, :] - X[:, 0, :]
            self.X = torch.FloatTensor(delta[:, np.newaxis, :])
        else:
            raise ValueError(f"Unknown mode: {mode}")
 
    def __len__(self):
        return len(self.y)
 
    def __getitem__(self, idx):
        return self.X[idx], self.y[idx]
 
print("Dataset class defined ✓")
\end{lstlisting}
 
\subsubsection*{8. Model Architecture}
 
\begin{lstlisting}
class SimpleTemporalNetwork(nn.Module):
    """
    Simplified model for small datasets.
    """
 
    def __init__(self, n_features, n_timepoints=2, hidden_dim=32, n_classes=3, dropout=0.5):
        super().__init__()
 
        # Simple 1D Conv
        self.conv = nn.Sequential(
            nn.Conv1d(n_timepoints, hidden_dim, kernel_size=5, padding=2),
            nn.BatchNorm1d(hidden_dim),
            nn.ReLU(),
            nn.Dropout(dropout)
        )
 
        # Simple LSTM
        self.lstm = nn.LSTM(
            input_size=hidden_dim,
            hidden_size=hidden_dim // 2,
            num_layers=1,
            batch_first=True,
            bidirectional=True
        )
 
        # Attention
        self.attention = nn.Sequential(
            nn.Linear(hidden_dim, hidden_dim // 2),
            nn.Tanh(),
            nn.Linear(hidden_dim // 2, 1)
        )
 
        # Classifier
        self.classifier = nn.Sequential(
            nn.Dropout(dropout),
            nn.Linear(hidden_dim, n_classes)
        )
 
    def forward(self, x):
        # CNN
        conv_out = self.conv(x)
 
        # LSTM
        lstm_in = conv_out.transpose(1, 2)
        lstm_out, _ = self.lstm(lstm_in)
 
        # Attention
        attn_scores = self.attention(lstm_out)
        attn_weights = F.softmax(attn_scores, dim=1)
        context = (lstm_out * attn_weights).sum(dim=1)
 
        # Classify
        logits = self.classifier(context)
 
        return logits
 
 
# Test model
n_features = X.shape[2]
test_model = SimpleTemporalNetwork(n_features=n_features, n_classes=len(class_names))
 
n_params = sum(p.numel() for p in test_model.parameters())
print(f"Model parameters: {n_params:,}")
print(f"Samples: {len(y)}")
print(f"Ratio (samples/params): {len(y)/n_params:.2f}")
 
# Test forward
dummy = torch.randn(4, 2, n_features)
out = test_model(dummy)
print(f"Forward pass OK: {dummy.shape} → {out.shape}")
\end{lstlisting}
 
\subsubsection*{9. Training Functions}
 
\begin{lstlisting}
def train_epoch(model, loader, optimizer, criterion, device):
    model.train()
    total_loss = 0
    correct = 0
    total = 0
 
    for X_batch, y_batch in loader:
        X_batch, y_batch = X_batch.to(device), y_batch.to(device)
 
        optimizer.zero_grad()
        outputs = model(X_batch)
        loss = criterion(outputs, y_batch)
 
        if torch.isnan(loss):
            continue
 
        loss.backward()
        torch.nn.utils.clip_grad_norm_(model.parameters(), max_norm=1.0)
        optimizer.step()
 
        total_loss += loss.item()
        _, predicted = outputs.max(1)
        correct += (predicted == y_batch).sum().item()
        total += y_batch.size(0)
 
    avg_loss = total_loss / len(loader) if len(loader) > 0 else 0
    accuracy = correct / total if total > 0 else 0
    return avg_loss, accuracy
 
 
def evaluate(model, loader, criterion, device):
    model.eval()
    total_loss = 0
    all_preds = []
    all_targets = []
 
    with torch.no_grad():
        for X_batch, y_batch in loader:
            X_batch, y_batch = X_batch.to(device), y_batch.to(device)
            outputs = model(X_batch)
            loss = criterion(outputs, y_batch)
 
            if not torch.isnan(loss):
                total_loss += loss.item()
 
            _, predicted = outputs.max(1)
            all_preds.extend(predicted.cpu().numpy())
            all_targets.extend(y_batch.cpu().numpy())
 
    avg_loss = total_loss / len(loader) if len(loader) > 0 else 0
    accuracy = accuracy_score(all_targets, all_preds)
    return avg_loss, accuracy, np.array(all_preds), np.array(all_targets)
 
print("Training functions defined ✓")
\end{lstlisting}
 
\subsubsection*{10. Cross-Validation Function}
 
\begin{lstlisting}
def run_cross_validation(X, y, mode, config, n_folds=5):
    """
    Cross-validation για robust evaluation.
    """
    print(f"\n{'='*60}")
    print(f"EXPERIMENT: {mode.upper()} MODEL ({n_folds}-fold CV)")
    print(f"{'='*60}")
 
    n_timepoints = 2 if mode == 'temporal' else 1
    n_features = X.shape[2]
    n_classes = len(np.unique(y))
 
    # Class weights
    class_counts = np.bincount(y)
    class_weights = torch.FloatTensor(len(y) / (len(class_counts) * class_counts)).to(config.DEVICE)
 
    # Cross-validation
    kfold = StratifiedKFold(n_splits=n_folds, shuffle=True, random_state=config.RANDOM_SEED)
 
    fold_results = []
    all_preds = []
    all_targets = []
 
    for fold, (train_idx, val_idx) in enumerate(kfold.split(X, y)):
        print(f"\nFold {fold+1}/{n_folds}", end=" ")
 
        # Split
        X_train, X_val = X[train_idx], X[val_idx]
        y_train, y_val = y[train_idx], y[val_idx]
 
        # Datasets
        train_dataset = ProteinDataset(X_train, y_train, mode=mode)
        val_dataset = ProteinDataset(X_val, y_val, mode=mode)
 
        train_loader = DataLoader(train_dataset, batch_size=config.BATCH_SIZE, shuffle=True)
        val_loader = DataLoader(val_dataset, batch_size=config.BATCH_SIZE)
 
        # Model
        model = SimpleTemporalNetwork(
            n_features=n_features,
            n_timepoints=n_timepoints,
            hidden_dim=config.HIDDEN_DIM,
            n_classes=n_classes,
            dropout=config.DROPOUT
        ).to(config.DEVICE)
 
        # Training setup
        criterion = nn.CrossEntropyLoss(weight=class_weights)
        optimizer = torch.optim.AdamW(
            model.parameters(),
            lr=config.LEARNING_RATE,
            weight_decay=config.WEIGHT_DECAY
        )
        scheduler = torch.optim.lr_scheduler.ReduceLROnPlateau(
            optimizer, mode='min', factor=0.5, patience=5
        )
 
        # Training loop
        best_val_acc = 0
        best_model_state = None
        patience_counter = 0
 
        for epoch in range(config.EPOCHS):
            train_loss, train_acc = train_epoch(model, train_loader, optimizer, criterion, config.DEVICE)
            val_loss, val_acc, _, _ = evaluate(model, val_loader, criterion, config.DEVICE)
 
            scheduler.step(val_loss)
 
            if val_acc > best_val_acc:
                best_val_acc = val_acc
                best_model_state = {k: v.cpu().clone() for k, v in model.state_dict().items()}
                patience_counter = 0
            else:
                patience_counter += 1
 
            if patience_counter >= config.PATIENCE:
                break
 
        # Evaluate best model
        if best_model_state is not None:
            model.load_state_dict(best_model_state)
        _, val_acc, preds, targets = evaluate(model, val_loader, criterion, config.DEVICE)
 
        fold_results.append(val_acc)
        all_preds.extend(preds)
        all_targets.extend(targets)
 
        print(f"→ Accuracy: {val_acc:.4f}")
 
    # Summary
    mean_acc = np.mean(fold_results)
    std_acc = np.std(fold_results)
 
    print(f"\n" + "-"*40)
    print(f"Mean Accuracy: {mean_acc:.4f} ± {std_acc:.4f}")
 
    return {
        'mean_acc': mean_acc,
        'std_acc': std_acc,
        'fold_results': fold_results,
        'all_preds': np.array(all_preds),
        'all_targets': np.array(all_targets)
    }
\end{lstlisting}
 
\subsubsection*{11. Run All Experiments ��}
 
\begin{lstlisting}
# Run all three experiments
results = {}
 
for mode in ['static', 'delta', 'temporal']:
    results[mode] = run_cross_validation(X, y, mode, Config, n_folds=Config.N_FOLDS)
\end{lstlisting}
 
\subsubsection*{12. Results Comparison ��}
 
\begin{lstlisting}
# Summary table
print("\n" + "="*60)
print("FINAL COMPARISON (5-Fold Cross-Validation)")
print("="*60)
print(f"\n{'Model':<15} {'Mean Accuracy':<18} {'Std':<10}")
print("-" * 45)
 
for mode in ['static', 'delta', 'temporal']:
    mean = results[mode]['mean_acc']
    std = results[mode]['std_acc']
    print(f"{mode:<15} {mean:<18.4f} ±{std:<10.4f}")
 
# Baseline comparison
random_baseline = 1.0 / len(class_names)
print(f"\n{'random guess':<15} {random_baseline:<18.4f}")
 
# Conclusion
print("\n" + "-"*45)
temporal_acc = results['temporal']['mean_acc']
static_acc = results['static']['mean_acc']
delta_acc = results['delta']['mean_acc']
 
best_mode = max(['static', 'delta', 'temporal'], key=lambda m: results[m]['mean_acc'])
print(f"Best model: {best_mode.upper()} ({results[best_mode]['mean_acc']:.4f})")
 
if temporal_acc > static_acc + 0.02:
    improvement = (temporal_acc - static_acc) * 100
    print(f"\n✓ Temporal > Static by {improvement:.1f}%")
    print("  → Evidence that temporal info helps!")
elif static_acc > temporal_acc + 0.02:
    print("\n⚠ Static performs better than Temporal")
    print("  → May need more data for temporal patterns")
else:
    print("\n~ Results are similar (within noise)")
    print("  → Need more data for definitive conclusions")
\end{lstlisting}
 
\begin{lstlisting}
# Visualization
fig, axes = plt.subplots(1, 2, figsize=(14, 5))
 
# Bar plot comparison
modes = ['static', 'delta', 'temporal']
means = [results[m]['mean_acc'] for m in modes]
stds = [results[m]['std_acc'] for m in modes]
 
colors = ['#3498db', '#e74c3c', '#2ecc71']
bars = axes[0].bar(modes, means, yerr=stds, capsize=5, color=colors, alpha=0.8)
axes[0].set_ylabel('Accuracy', fontsize=12)
axes[0].set_title('Model Comparison (5-Fold CV)', fontsize=14)
axes[0].set_ylim(0, 1)
axes[0].axhline(y=1/len(class_names), color='gray', linestyle='--', label='Random guess')
axes[0].legend()
axes[0].grid(axis='y', alpha=0.3)
 
# Add values on bars
for bar, mean, std in zip(bars, means, stds):
    axes[0].text(bar.get_x() + bar.get_width()/2, mean + std + 0.03,
                 f'{mean:.3f}', ha='center', fontsize=11, fontweight='bold')
 
# Confusion matrix for best model
best_mode = max(modes, key=lambda m: results[m]['mean_acc'])
cm = confusion_matrix(results[best_mode]['all_targets'], results[best_mode]['all_preds'])
 
sns.heatmap(cm, annot=True, fmt='d', cmap='Blues',
            xticklabels=class_names, yticklabels=class_names, ax=axes[1],
            annot_kws={"size": 14})
axes[1].set_xlabel('Predicted', fontsize=12)
axes[1].set_ylabel('True', fontsize=12)
axes[1].set_title(f'Confusion Matrix ({best_mode.capitalize()} Model)', fontsize=14)
 
plt.tight_layout()
plt.show()
\end{lstlisting}
 
\begin{lstlisting}
# Classification report for best model
best_mode = max(modes, key=lambda m: results[m]['mean_acc'])
 
print(f"\nClassification Report ({best_mode.capitalize()} Model):")
print("="*50)
print(classification_report(
    results[best_mode]['all_targets'],
    results[best_mode]['all_preds'],
    target_names=class_names
))
\end{lstlisting}
 
\subsubsection*{13. Top Discriminative Proteins}
 
\begin{lstlisting}
# Get protein importance scores
scores = feature_selector.scores_[selected_indices]
protein_importance = sorted(zip(selected_proteins, scores), key=lambda x: x[1], reverse=True)
 
print("\nTop 20 Most Discriminative Proteins:")
print("="*45)
for i, (protein, score) in enumerate(protein_importance[:20]):
    print(f"{i+1:2d}. {protein:<20} (F-score: {score:.2f})")
\end{lstlisting}
 
\begin{lstlisting}
# Visualize top proteins
top_n = 15
top_proteins = [p for p, _ in protein_importance[:top_n]]
top_scores = [s for _, s in protein_importance[:top_n]]
 
fig, ax = plt.subplots(figsize=(10, 6))
bars = ax.barh(range(top_n), top_scores[::-1], color='steelblue', alpha=0.8)
ax.set_yticks(range(top_n))
ax.set_yticklabels(top_proteins[::-1])
ax.set_xlabel('F-score (ANOVA)', fontsize=12)
ax.set_title('Top 15 Discriminative Proteins', fontsize=14)
ax.grid(axis='x', alpha=0.3)
 
plt.tight_layout()
plt.show()
\end{lstlisting}
 
\subsubsection*{14. Save Results}
 
\begin{lstlisting}
# Save results to CSV
results_df = pd.DataFrame({
    'Model': ['Static', 'Delta', 'Temporal'],
    'Mean_Accuracy': [results['static']['mean_acc'], results['delta']['mean_acc'], results['temporal']['mean_acc']],
    'Std_Accuracy': [results['static']['std_acc'], results['delta']['std_acc'], results['temporal']['std_acc']],
    'Fold_1': [results['static']['fold_results'][0], results['delta']['fold_results'][0], results['temporal']['fold_results'][0]],
    'Fold_2': [results['static']['fold_results'][1], results['delta']['fold_results'][1], results['temporal']['fold_results'][1]],
    'Fold_3': [results['static']['fold_results'][2], results['delta']['fold_results'][2], results['temporal']['fold_results'][2]],
    'Fold_4': [results['static']['fold_results'][3], results['delta']['fold_results'][3], results['temporal']['fold_results'][3]],
    'Fold_5': [results['static']['fold_results'][4], results['delta']['fold_results'][4], results['temporal']['fold_results'][4]],
})
 
results_df.to_csv('experiment_results.csv', index=False)
print("Results saved to 'experiment_results.csv'")
display(results_df)
\end{lstlisting}
 
\begin{lstlisting}
# Save top proteins
proteins_df = pd.DataFrame({
    'Rank': range(1, len(protein_importance) + 1),
    'Protein': [p for p, _ in protein_importance],
    'F_score': [s for _, s in protein_importance]
})
 
proteins_df.to_csv('top_proteins.csv', index=False)
print("Top proteins saved to 'top_proteins.csv'")
display(proteins_df.head(20))
\end{lstlisting}
 
\subsubsection*{Summary}
 
